\documentclass[11pt,a4paper]{article}

\usepackage[utf8]{inputenc}
\usepackage[T1]{fontenc}
\usepackage[polish]{babel}
\usepackage{geometry}
\usepackage{booktabs}
\usepackage{tabularray}
\usepackage{afterpage}
\usepackage{rotating}
\usepackage{longtable}
\usepackage{array}
\usepackage{xcolor}
\usepackage{hyperref}
\usepackage{titlesec}
\usepackage{parskip}
\usepackage{enumitem}
\usepackage{fancyhdr}
\usepackage{pdflscape}
\usepackage{makecell}

\geometry{margin=2.5cm}

\definecolor{accent}{RGB}{30, 60, 100}
\definecolor{placeholdercolor}{RGB}{160, 160, 160}
\definecolor{rowA}{RGB}{235, 242, 250}

\titleformat{\section}{\large\bfseries\color{accent}}{}{0em}{}[\titlerule]
\titleformat{\subsection}{\normalsize\bfseries}{}{0em}{}
\titleformat{\subsubsection}{\normalsize\itshape}{}{0em}{}

\pagestyle{fancy}
\fancyhf{}
\rhead{\small Opis Ról i Obowiązków}
\lhead{\small Projekt: \textit{Feedback}}
\cfoot{\thepage}
\renewcommand{\headrulewidth}{0.4pt}

\hypersetup{colorlinks=true, linkcolor=accent, urlcolor=accent}

% ── MAKRA RACI ───────────────────────────────────────────────────
% ph = placeholder (do ustalenia)
\newcommand{\ph}{{\color{placeholdercolor}?}}
% Skróty ról RACI
\newcommand{\R}{\textbf{R}}
\newcommand{\Ac}{\textbf{A}}
\newcommand{\C}{\textbf{C}}
\newcommand{\I}{\textbf{I}}
\newcommand{\Su}{\textbf{S}}
\newcommand{\none}{--}

% ── PRZYPISANIA RACI ─────────────────────────────────────────────
% Wiersze: zadania | Kolumny: KP, AB, RG, JP, MP, FS, KS, GD, GU
% KP = Trzciński, AB = Bielicki, RG = Głasek, JP = Pączka,
% MP = Popkowicz, FS = Shramko
% KS = Sobótka (Komitet Sterujący), GD = Płocharski (Gł. Dostawca), GU = Herman (Gł. Użytkownik)

% 1. Dokumentacja PRINCE2
\newcommand{\raciDocKP}{\R}
\newcommand{\raciDocAB}{\Su}
\newcommand{\raciDocRG}{\Su}
\newcommand{\raciDocJP}{\Su}
\newcommand{\raciDocMP}{\Su}
\newcommand{\raciDocFS}{\Su}
\newcommand{\raciDocKS}{\Ac}
\newcommand{\raciDocGD}{\C}
\newcommand{\raciDocGU}{\I}

% 2. Mechanika parowania (backend core)
\newcommand{\raciMechKP}{\Ac}
\newcommand{\raciMechAB}{\ph}
\newcommand{\raciMechRG}{\ph}
\newcommand{\raciMechJP}{\ph}
\newcommand{\raciMechMP}{\ph}
\newcommand{\raciMechFS}{\ph}
\newcommand{\raciMechKS}{\I}
\newcommand{\raciMechGD}{\C}
\newcommand{\raciMechGU}{\I}

% 3. Frontend / UI / HUD
\newcommand{\raciFrontKP}{\Ac}
\newcommand{\raciFrontAB}{\ph}
\newcommand{\raciFrontRG}{\ph}
\newcommand{\raciFrontJP}{\ph}
\newcommand{\raciFrontMP}{\ph}
\newcommand{\raciFrontFS}{\ph}
\newcommand{\raciFrontKS}{\I}
\newcommand{\raciFrontGD}{\C}
\newcommand{\raciFrontGU}{\C}

% 4. Projektowanie poziomów i zagadek
\newcommand{\raciLvlKP}{\Ac}
\newcommand{\raciLvlAB}{\ph}
\newcommand{\raciLvlRG}{\ph}
\newcommand{\raciLvlJP}{\ph}
\newcommand{\raciLvlMP}{\ph}
\newcommand{\raciLvlFS}{\ph}
\newcommand{\raciLvlKS}{\I}
\newcommand{\raciLvlGD}{\C}
\newcommand{\raciLvlGU}{\C}

% 5. Grafika i assety
\newcommand{\raciGrafKP}{\Ac}
\newcommand{\raciGrafAB}{\ph}
\newcommand{\raciGrafRG}{\ph}
\newcommand{\raciGrafJP}{\ph}
\newcommand{\raciGrafMP}{\ph}
\newcommand{\raciGrafFS}{\ph}
\newcommand{\raciGrafKS}{\I}
\newcommand{\raciGrafGD}{\none}
\newcommand{\raciGrafGU}{\C}

% 6. Testy i kontrola jakości
\newcommand{\raciTestKP}{\Ac}
\newcommand{\raciTestAB}{\ph}
\newcommand{\raciTestRG}{\ph}
\newcommand{\raciTestJP}{\ph}
\newcommand{\raciTestMP}{\ph}
\newcommand{\raciTestFS}{\ph}
\newcommand{\raciTestKS}{\I}
\newcommand{\raciTestGD}{\Ac}
\newcommand{\raciTestGU}{\C}

% 7. Code review i zarządzanie repozytorium
\newcommand{\raciCRKP}{\Ac}
\newcommand{\raciCRAB}{\ph}
\newcommand{\raciCRRG}{\ph}
\newcommand{\raciCRJP}{\ph}
\newcommand{\raciCRMP}{\ph}
\newcommand{\raciCRFS}{\ph}
\newcommand{\raciCRKS}{\none}
\newcommand{\raciCRGD}{\none}
\newcommand{\raciCRGU}{\none}

% 8. Prezentacja finalna i demo
\newcommand{\raciPrezKP}{\R}
\newcommand{\raciPrezAB}{\Su}
\newcommand{\raciPrezRG}{\Su}
\newcommand{\raciPrezJP}{\Su}
\newcommand{\raciPrezMP}{\Su}
\newcommand{\raciPrezFS}{\Su}
\newcommand{\raciPrezKS}{\Ac}
\newcommand{\raciPrezGD}{\Ac}
\newcommand{\raciPrezGU}{\C}

% ─────────────────────────────────────────────────────────────────

\begin{document}

% ── STRONA TYTUŁOWA ──────────────────────────────────────────────
\begin{titlepage}
\centering
\vspace*{2cm}
{\LARGE\bfseries\color{accent} Opis Ról i Obowiązków}\\[0.5cm]
{\large Projekt: \textit{Feedback}}\\[0.3cm]
{\normalsize Mechanika parowania jako klucz do sukcesu}\\[2cm]

\begin{tabular}{ll}
\textbf{Kierownik projektu:}   & Kacper Trzciński \\[0.3em]
\textbf{Przewodnicząca KS:}    & Aleksandra Sobótka \\[0.3em]
\textbf{Główny Dostawca:}      & Aleksander Płocharski \\[0.3em]
\textbf{Główny Użytkownik:}    & Tomasz Herman \\[0.3em]
\textbf{Zespół:}               & A. Bielicki, R. Głasek, J. Pączka, \\
                               & M. Popkowicz, F. Shramko, K. Trzciński \\[0.3em]
\textbf{Data opracowania:}     & \today \\[0.3em]
\end{tabular}

\vfill
{\small Zarządzanie Przedsięwzięciami Informatycznymi\\
Wydział Matematyki i Nauk Informacyjnych, Politechnika Warszawska}
\end{titlepage}

% ── SPIS TREŚCI ─────────────────────────────────────────────────
\tableofcontents
\newpage

% ════════════════════════════════════════════════════════════════
\section{Wstęp}
% ════════════════════════════════════════════════════════════════

\subsection{Przeznaczenie dokumentu}
Niniejszy dokument opisuje strukturę organizacyjną projektu \textit{Feedback} oraz definiuje role, obowiązki i zakres odpowiedzialności wszystkich jego uczestników. Dokument stanowi uzupełnienie Karty projektu.

\subsection{Po co definiujemy role i obowiązki?}
W projekcie realizowanym przez 6-osobowy zespół studencki, działający równolegle w kilku obszarach (implementacja, grafika, game design, testy), szczególnie łatwo o sytuację, w której zadanie wypada między odpowiedzialnościami dwóch osób i nie zostaje wykonane przez nikogo. Precyzyjne zdefiniowanie ról ma na celu:

\begin{itemize}[noitemsep]
    \item jednoznaczne określenie, kto podejmuje decyzje w każdym obszarze projektu;
    \item uniknięcie dublowania pracy lub powstawania luk odpowiedzialności;
    \item ustalenie, do kogo należy się zwrócić po konsultację lub zatwierdzenie;
    \item zapewnienie ciągłości projektu w przypadku zmiany składu zespołu.
\end{itemize}

% ════════════════════════════════════════════════════════════════
\section{Struktura organizacyjna projektu wg PRINCE2}
% ════════════════════════════════════════════════════════════════

PRINCE2 definiuje trójwarstwową strukturę zarządzania projektem: Komitet Sterujący (poziom strategiczny), Kierownik projektu (poziom zarządczy) oraz Zespół projektowy (poziom wytwórczy). Poniżej opisano obsadę i obowiązki każdej wymaganej przez metodykę roli.

\subsection{Komitet Sterujący}

Komitet Sterujący (KS) ponosi ogólną odpowiedzialność za powodzenie projektu. Zatwierdza kluczowe decyzje, monitoruje zgodność z celami biznesowymi i reaguje na eskalowane ryzyka. W skład KS wchodzą trzy role:

\subsubsection{Przewodniczący Komitetu Sterującego}
\textbf{Obsada:} Aleksandra Sobótka (prowadząca przedmiot w semestrze dokumentacyjnym)

\textbf{Obowiązki:}
\begin{itemize}[noitemsep]
    \item zatwierdzanie kluczowych dokumentów projektowych (Karta projektu, WBS, Rejestr ryzyk);
    \item ocena postępów na kolejnych etapach semestrów;
    \item rozstrzyganie eskalowanych zagadnień i decyzji wykraczających poza kompetencje kierownika projektu;
    \item zatwierdzanie zmian zakresu projektu.
\end{itemize}

\subsubsection{Główny Użytkownik}
\textbf{Obsada:} Tomasz Herman (asystent badawczo-dydaktyczny, regularnie uczestniczący w zajęciach)

\textbf{Obowiązki:}
\begin{itemize}[noitemsep]
    \item reprezentowanie perspektywy końcowego odbiorcy gry (gracza);
    \item weryfikacja, czy wymagania dotyczące rozgrywki i mechanik są spełnione;
    \item uczestnictwo w playtestach i przekazywanie feedbacku na temat grywalności i czytelności mechaniki parowania;
    \item zatwierdzenie finalnej wersji gry jako spełniającej oczekiwania użytkownika.
\end{itemize}

\subsubsection{Główny Dostawca}
\textbf{Obsada:} Aleksander Płocharski (prowadzący przedmiot w semestrze implementacyjnym)

\textbf{Obowiązki:}
\begin{itemize}[noitemsep]
    \item nadzór nad jakością techniczną dostarczanego produktu w semestrze letnim 2026/2027;
    \item konsultacje dotyczące wyboru technologii, architektury i rozwiązań implementacyjnych;
    \item zatwierdzanie wyników testów jako potwierdzenie gotowości do prezentacji finalnej;
    \item weryfikacja, czy zespół dysponuje odpowiednimi zasobami i kompetencjami do realizacji zakresu.
\end{itemize}

% \noindent\textit{Uwaga: projekt realizowany jest dwuetapowo. W semestrze dokumentacyjnym (zimowym 2025/2026) rolę wiodącą w KS pełni Aleksandra Sobótka. Aleksander Płocharski obejmuje aktywny nadzór nad projektem od początku semestru letniego 2026/2027, kiedy rozpoczyna się faza implementacji.}

\subsection{Kierownik Projektu}

\textbf{Obsada:} Kacper Trzciński

\textbf{Obowiązki:}
\begin{itemize}[noitemsep]
    \item planowanie i koordynowanie pracy zespołu zgodnie z harmonogramem i metodyką PRINCE2;
    \item prowadzenie codziennej komunikacji wewnętrznej oraz jedyna formalna komunikacja z KS i prowadzącymi;
    \item prowadzenie i aktualizowanie rejestru ryzyk, rejestru zagadnień i harmonogramu;
    \item sporządzanie raportów końca sprintu;
    \item zarządzanie zmianami zakresu i eskalowanie decyzji do KS;
    \item prowadzenie i archiwizowanie notatek ze spotkań.
\end{itemize}

\textit{W projekcie tej skali Kierownik Projektu pełni jednocześnie role Kierownika Zespołu, Nadzoru Projektu oraz Wsparcia Projektu, zgodnie z dopuszczeniem PRINCE2 dla małych projektów.}

\subsection{Zespół Projektowy}

Sześciu członków zespołu realizuje zadania wytwórcze. Szczegółowy podział kompetencji i odpowiedzialności opisano w sekcji~\ref{sec:zespol} oraz w macierzy RACI (sekcja~\ref{sec:raci}).

% ════════════════════════════════════════════════════════════════
\section{Zespół projektowy --- kompetencje i obszary}
\label{sec:zespol}
% ════════════════════════════════════════════════════════════════

\subsection{Wymagane kompetencje}

Do realizacji projektu \textit{Feedback} niezbędne są następujące kompetencje:

\begin{description}[style=nextline]
    \item[Programowanie (backend gry)]
    Implementacja mechaniki parowania (zmiana momentum, odbijanie pocisków), systemu kolizji i fizyki, logiki przeciwników, obsługi stanu gry. Wymagana znajomość wybranego silnika gry (Godot / Unity / Unreal) oraz podstaw fizyki 2D.

    \item[Programowanie (frontend / UI)]
    Implementacja interfejsu użytkownika (menu, HUD, ekran końcowy), obsługa sterowania postacią, integracja animacji i przejść między ekranami.

    \item[Game design / projektowanie poziomów]
    Projektowanie ekranów z zagadkami logicznymi opartymi na mechanice parowania, definiowanie typów przeciwników i ich roli w zagadkach, pacing i balans rozgrywki.

    \item[Grafika 2D / UI art]
    Tworzenie sprite'ów postaci, przeciwników i elementów poziomów, animacje ataku i parowania, projekt elementów interfejsu.

    \item[Testowanie]
    Weryfikacja poprawności mechaniki parowania, testowanie poziomów logicznych, testy systemu kolizji i wydajności, testy końcowe wersji grywalnej, rejestrowanie i priorytetyzacja błędów.

    \item[Dokumentacja i zarządzanie projektem]
    Przygotowywanie dokumentacji PRINCE2, prowadzenie harmonogramu (ProjectLibre), sporządzanie raportów i notatek ze spotkań. Kompetencja obsadzana przez Kierownika Projektu.
\end{description}

\subsection{Skład zespołu i planowane obszary odpowiedzialności}

\begin{longtblr}[
    caption = {Planowany podział obszarów w zespole},
    label   = {tab:zespol},
]{
    colspec = {p{4cm} p{4cm} p{7cm}},
    rowhead = 1,
    hlines,
    row{1} = {font=\bfseries},
}
Imię i nazwisko & Rola projektowa & Planowany obszar odpowiedzialności \\
Kacper Trzciński    & Kierownik projektu & Koordynacja, dokumentacja PRINCE2, komunikacja z KS \\
Aleksander Bielicki & \ph               & \ph \\
Radosław Głasek     & \ph               & \ph \\
Jakub Pączka        & \ph               & \ph \\
Michał Popkowicz    & \ph               & \ph \\
Filip Shramko       & \ph               & \ph \\
\end{longtblr}

\noindent\textit{Szczegółowy podział ról zostanie ustalony na początku fazy implementacji (sprint 0) na podstawie zdeklarowanych kompetencji i preferencji członków zespołu.}

% ════════════════════════════════════════════════════════════════
\section{Macierz RACI}
\label{sec:raci}
% ════════════════════════════════════════════════════════════════

\subsection{Legenda}

Macierz RACI (w rozszerzeniu RASCI) przypisuje każdemu zadaniu projektowemu jedną z następujących ról dla każdego uczestnika:

\begin{description}[style=nextline, noitemsep]
    \item[\textbf{R} --- Responsible (Odpowiedzialny wykonawczo)]
    Osoba lub osoby, które faktycznie wykonują dane zadanie. Dbają o to, żeby praca została zrobiona. Za każde zadanie powinna być wyznaczona co najmniej jedna osoba z~rolą~R.
    \item[\textbf{A} --- Accountable (Odpowiedzialny decyzyjnie)]
    Jedna osoba, która zatwierdza wynik i ponosi ostateczną odpowiedzialność za jakość. Tylko jedna osoba na zadanie.
    \item[\textbf{S} --- Support (Wsparcie)]
    Osoby aktywnie uczestniczące w wykonaniu zadania, choć nie prowadzące go. Pomagają w realizacji, gdy są o to poproszone.
    \item[\textbf{C} --- Consult (Konsultowany)]
    Osoby, które nie wykonują zadania, ale dostarczają wiedzy lub informacji niezbędnych do jego realizacji. Komunikacja dwukierunkowa.
    \item[\textbf{I} --- Inform (Informowany)]
    Osoby, które otrzymują informację o postępie lub wynikach zadania, ale nie są zaangażowane w jego realizację. Komunikacja jednostronna.
\end{description}

\subsection{Skróty kolumn}

\begin{tabular}{ll@{\hspace{2em}}ll}
\textbf{KP} & Kacper Trzciński (Kierownik Projektu) &
\textbf{MP} & Michał Popkowicz \\
\textbf{AB} & Aleksander Bielicki &
\textbf{FS} & Filip Shramko \\
\textbf{RG} & Radosław Głasek &
\textbf{KS} & Aleksandra Sobótka (Komitet Sterujący) \\
\textbf{JP} & Jakub Pączka &
\textbf{GD} & Aleksander Płocharski (Główny Dostawca) \\
   &  &
\textbf{GU} & Tomasz Herman (Główny Użytkownik) \\
\end{tabular}

\vspace{1em}
\noindent Komórki oznaczone {\color{placeholdercolor}\textbf{?}} wymagają ustalenia na etapie planowania sprintu~0.

\vspace{1em}
\noindent\textbf{Przegląd dokumentu:} macierz RACI zostanie uzupełniona i zatwierdzona przez Kierownika Projektu na początku fazy implementacyjnej (sprint~0), po ustaleniu obszarów odpowiedzialności poszczególnych członków zespołu. Dokument podlega aktualizacji przy każdej istotnej zmianie składu lub zakresu projektu.

\clearpage
\begin{landscape}

\begin{longtblr}[
    caption = {Macierz RACI projektu Feedback},
    label   = {tab:raci},
]{
    colspec = {p{5.5cm} *{9}{X[c]}},
    rowhead = 1,
    hlines,
    row{1} = {font=\bfseries, bg=accent, fg=white},
    row{3,5,7,9} = {bg=rowA},
}
Zadanie
    & KP & AB & RG & JP & MP & FS & KS & GD & GU \\
Dokumentacja PRINCE2
    & \raciDocKP & \raciDocAB & \raciDocRG & \raciDocJP
    & \raciDocMP & \raciDocFS & \raciDocKS & \raciDocGD & \raciDocGU \\
Implementacja mechaniki parowania (backend)
    & \raciMechKP & \raciMechAB & \raciMechRG & \raciMechJP
    & \raciMechMP & \raciMechFS & \raciMechKS & \raciMechGD & \raciMechGU \\
Frontend -- UI, HUD, sterowanie
    & \raciFrontKP & \raciFrontAB & \raciFrontRG & \raciFrontJP
    & \raciFrontMP & \raciFrontFS & \raciFrontKS & \raciFrontGD & \raciFrontGU \\
Projektowanie poziomów i zagadek
    & \raciLvlKP & \raciLvlAB & \raciLvlRG & \raciLvlJP
    & \raciLvlMP & \raciLvlFS & \raciLvlKS & \raciLvlGD & \raciLvlGU \\
Grafika i assety (sprite'y, animacje)
    & \raciGrafKP & \raciGrafAB & \raciGrafRG & \raciGrafJP
    & \raciGrafMP & \raciGrafFS & \raciGrafKS & \raciGrafGD & \raciGrafGU \\
Testy i kontrola jakości
    & \raciTestKP & \raciTestAB & \raciTestRG & \raciTestJP
    & \raciTestMP & \raciTestFS & \raciTestKS & \raciTestGD & \raciTestGU \\
Code review i zarządzanie repozytorium
    & \raciCRKP & \raciCRAB & \raciCRRG & \raciCRJP
    & \raciCRMP & \raciCRFS & \raciCRKS & \raciCRGD & \raciCRGU \\
Prezentacja finalna i demo gry
    & \raciPrezKP & \raciPrezAB & \raciPrezRG & \raciPrezJP
    & \raciPrezMP & \raciPrezFS & \raciPrezKS & \raciPrezGD & \raciPrezGU \\
\end{longtblr}

\end{landscape}
\end{document}